% Methodology
\section{Methodology}

\subsection{Data Sources}

Our analysis is based on the corpus of approximately 4,000 known Indus inscriptions, using:

\begin{itemize}
    \item Mahadevan's concordance (1977) for sign classification
    \item Wells' supplementary concordance for variants
    \item Site distribution data from major excavations (Harappa, Mohenjo-daro, Dholavira, etc.)
    \item Object type associations (seals, tablets, pottery, etc.)
\end{itemize}

\subsection{Analytical Framework}

Our methodology combines:

\subsubsection{Distributional Analysis}
\begin{itemize}
    \item Frequency analysis of individual signs
    \item Positional distribution (initial, medial, terminal)
    \item Bigram and trigram analysis
    \item Site-specific and object-specific patterns
\end{itemize}

\subsubsection{Combinatorial Analysis}
\begin{itemize}
    \item Valid sequence constraints
    \item Sign co-occurrence restrictions
    \item Length distribution patterns
    \item Structural templates and schemas
\end{itemize}

\subsubsection{Formal Modeling}
\begin{itemize}
    \item Finite automaton construction
    \item State space analysis
    \item Transition function derivation
    \item Accepting condition specification
\end{itemize}

\subsubsection{Comparative Analysis}
\begin{itemize}
    \item Egyptian hieroglyphic determinatives
    \item Old Chinese oracle bone notation
    \item Phaistos Disk structure
    \item Modern operational systems (JSON, XML schemas, workflow notations)
\end{itemize}

\subsection{Validation Strategy}

Claims are validated through:
\begin{enumerate}
    \item Consistency with known corpus patterns
    \item Predictive power for unseen inscriptions
    \item Cross-validation against independent data subsets
    \item Comparison with functionally analogous systems
\end{enumerate}
